A managed run-time owes its users two numbers it rarely states: how much memory it will use, and how long it may pause.
Collectors answer neither, because a collector is a scheduler: it decides at run time when to work, how much to work, and how long the mutator waits.
This paper describes extC, a systems language whose reclamation is decided entirely at compile time and whose run-time work is a bounded, incremental algorithm, and it turns the two numbers into contracts that can be checked.

Three ideas carry the design.
First, the compiler places every allocation in one of five storage substrates --- frame arenas, typed reclamation regions, container pools, per-task arenas, and address-reserved plates --- so that four of the five have a scope-derived lifetime and need no collector at all.
Second, for the one substrate that can hold cycles the compiler emits one of three tiers per translation unit: no reclamation code at all (544 of the 794 corpus programs that compile, 68.5\%), an event tier that counts and never sweeps, or an incremental mark-and-sweep tier whose work is cut into segments by a wall-clock deadline, charged against an allocation quota, and whose region destruction is deferred through an $O(1)$ closing queue.
Third, every run-time decision that could grow without bound is given a budget and a proof: the pause of one step is bounded by a segment length, the heap is bounded by the live set through a trigger law, and the admission of an allocation reserves room for the largest request so that a single object can never overrun the wall.

We prove the safety of the scheme (no dangling reference) together with tier soundness, idle transparency, bounded memory, bounded pause, amortized work, and the progress of deferred destruction, and we machine-check the core of the model in Lean~4.
We then formalize reclamation policy as a roofline: the cost per allocated byte is $\lambda(r) = (c_t + c_s F r)/(r-1)$, memory follows $M = L + \kappa_M H$, a step of $b$ pops costs $P(b) = b\tau_t + W_s 64 \tau_s$, and keeping up with the mutator requires $r \ge 1 + gK/b$; the admissible policies form a band on the $(\lambda, a)$ plane whose walls are the memory budget and the pause budget, and the optimal policy sits on both.
All six constants are measured, and the model tracks measurement to within 2--10\%.

On four cross-language workloads extC matches the best Rust arm on promotion (39.0 vs 39.1~ms), beats the best C++ arm on size-ratio (16.0 vs 17.9~ms) and on a ring workload (15.0 vs 36.1~ms), and is within 1.1\% of C++ on a pointer-chasing graph at 13\% lower peak memory, while no workload shows a single pause above one millisecond and the graph shape has a 99.9th-percentile pause of $58\,_\mu$s.
Turning reclamation off makes the allocation-heavy shapes 1.5--3.7 times slower and 17--59 times larger in memory, which is the cost the design removes.
